\documentclass[11pt]{article}
\usepackage[margin=1in]{geometry}
\usepackage{amsmath,amssymb,amsthm}
\usepackage{booktabs}
\usepackage{listings}
\usepackage{xcolor}
\usepackage[numbers]{natbib}
\usepackage{hyperref}

\newtheorem{theorem}{Theorem}
\theoremstyle{remark}
\newtheorem{remark}{Remark}
\newtheorem{example}{Example}

\lstset{basicstyle=\ttfamily\small,columns=fullflexible,breaklines=true,
  frame=single,keywordstyle=\color{blue!60!black},commentstyle=\color{gray},
  morekeywords={def,theorem,lemma,namespace,end,example,by,fun,import,noncomputable,section,open,intro,have,rw,at,exact,axiom,let,refine,obtain,rintro,unfold,constructor},
  morecomment=[l]{--},
  literate={ℝ}{{$\mathbb{R}$}}1 {→}{{$\to$}}1 {∀}{{$\forall$}}1 {∃}{{$\exists$}}1
           {¬}{{$\neg$}}1 {≤}{{$\le$}}1 {≥}{{$\ge$}}1 {↔}{{$\leftrightarrow$}}1
           {∧}{{$\wedge$}}1 {δ}{{$\delta$}}1 {·}{{$\cdot$}}1 {−}{{$-$}}1 {⁻}{{$^{-}$}}1
           {₀}{{$_0$}}1 {⟨}{{$\langle$}}1 {⟩}{{$\rangle$}}1 {Δ}{{$\Delta$}}1}

\title{A Correction Note on Two Early Intelligence Bound Sources:\\
The Unrestricted Axiom Is False, and ``No Finite Compensation'' Proves Finite Compensation}
\author{Justin Hart\\Viridis LLC}
\date{19 November 2026}

\begin{document}
\maketitle

\begin{abstract}
Two early Viridis Lean sources of January 2026 contain formal text that does not match its
prose. First, the axiom \texttt{intelligence\_bound} states
$\dot I\le \frac{P}{k_B T\ln 2}\,D$ for every real $\dot I$ under elementary guards, and is false
as written: the instance $\dot I=1$, $P=D=0$, $k_B=T=1$ reduces it to $1\le0$. Second, a theorem
named \texttt{no\_finite\_compensation}, introduced as saying that no finite power increase can
compensate a loss of the diversity factor $D$, in fact proves that the finite increase
$P\,\Delta D/(D-\Delta D)$ restores the bound exactly. We give Lean~4 proofs of two sealed
statements. \texttt{P5.UnrestrictedBoundFalse} shows that the unrestricted axiom is false.
\texttt{P5.CompensationUnderCeiling} separates four facts about the capacity
$\mathrm{cap}(P,D,C)=PD/C$ that the historical prose runs together: (1)~after a partial loss
$0<\Delta D<D$ the power $PD/(D-\Delta D)$ restores the capacity and is the only power that does;
(2)~under a power ceiling $P\le P_{\max}$, restoration is possible exactly when
$\Delta D\le D(1-P/P_{\max})$; (3)~the compensating power grows without bound as
$\Delta D\to D^-$; and (4)~at $D=0$ no power restores a positive capacity. In economic terms,
along a level set of $PD$ the two inputs substitute for each other. $D$ is essential, since
nothing replaces it at zero, but it is not non-substitutable in the interior; what makes a loss of
$D$ uncompensable is a finite power ceiling. Both results are elementary corrections of formal
text. They say nothing about how any agent behaves, about any measured value of $D$, or about
ecology.
\end{abstract}

\section{Introduction}
A machine-checked development is only as reliable as the match between its formal text and what
is said about it. A review of the Viridis research corpus on 11~November~2026 found two early Lean
declarations whose formal content diverges from their surrounding prose. Project P5 checks both
with the Lean kernel. The first session (P5-T001) proposed exact statements, which were then
sealed. The second session (P5-T002) proved the first correction. This session (P5-T003) proves
the second.

\paragraph{Finding 1.} The historical \texttt{intelligence\_bound} axiom quantifies over every
real rate $\dot I$ and is therefore false (Section~\ref{sec:axiom}). An environment that assumes
it is inconsistent.

\paragraph{Finding 2.} The historical theorem \texttt{no\_finite\_compensation} is described as
showing that ``there is no finite power increase $\Delta P$ that can fully compensate'' a loss of
$D$. Its formal statement, however, is that a specific finite $\Delta P$ \emph{does} compensate
(Section~\ref{sec:comp}). The prose mixes several true facts: finite compensation in the interior,
a limit imposed by a power ceiling, divergence near complete loss, and impossibility at $D=0$.
These have different hypotheses and different consequences. We state and prove them separately.

The questions matter because these sources were presented as formal support for a conclusion that
$D$ cannot be traded off against power. The correct reading is narrower. Within the capacity
algebra, $D$ is \emph{essential} (zero $D$ gives zero capacity whatever the power), but a partial
loss is always compensable with enough power. Only a finite ceiling on power makes some partial
losses uncompensable. None of this is new mathematics. The aim is to keep the record accurate.

\section{Background and related work}
Landauer's principle~\citep{landauer1961} associates a minimal dissipation of $k_B T\ln 2$ with
each logically irreversible bit operation at temperature $T$. Dividing a power budget $P$ by this
energy gives an upper bound on the rate of irreversible bit operations. The early sources define,
over the reals,
\[
  \mathrm{landauerEnergy}(k_B,T)=k_B T\ln 2,\qquad
  \mathrm{maxBitRate}(P,k_B,T)=\frac{P}{k_B T\ln 2},
\]
and multiply the result by a diversity factor $D\in[0,1]$. Writing $C=k_BT\ln2$ for the cost per
bit, the resulting capacity is
\[
  \mathrm{cap}(P,D,C)=\frac{P\,D}{C}.
\]
For fixed $C$ this is a Cobb--Douglas form in $(P,D)$ with unit exponents. The distinction between
an input being \emph{essential} (output is zero without it) and being \emph{substitutable} (a loss
of it can be offset by more of another input) is standard in the economics of exhaustible
resources, where it was central to the analyses of
Dasgupta and Heal~\citep{dasguptaheal1974} and Solow~\citep{solow1974}. We borrow only the
vocabulary. Nothing here depends on those models.

We use Lean~4~\citep{lean4} and Mathlib~\citep{mathlib2020}. The historical files are quoted, not
imported. The sealed statements restate the needed definitions over Mathlib alone.

\section{Finding 1: the unrestricted axiom is false}\label{sec:axiom}
The project README quotes the declaration in \texttt{Gaia\_Intelligence\_Theorem.lean}
(SHA-256 prefix \texttt{37c017145dcf\ldots2098c4}) as follows:
\begin{lstlisting}
axiom intelligence_bound :
  ∀ (dI_dt P D k_B T : ℝ),
  k_B > 0 → T > 0 → P ≥ 0 → D ≥ 0 → D ≤ 1 →
  dI_dt ≤ max_bit_rate P k_B T * D
\end{lstlisting}
The rate \texttt{dI\_dt} is a bound variable that ranges over every real number. Nothing in the
statement says whose rate it is.

\begin{theorem}[sealed statement \texttt{P5.UnrestrictedBoundFalse}]\label{thm:false}
It is not the case that for all real $\dot I,P,D,k_B,T$ with $k_B>0$, $T>0$, $P\ge0$, $D\ge0$
and $D\le1$ we have $\dot I\le\mathrm{maxBitRate}(P,k_B,T)\cdot D$.
\end{theorem}
\begin{proof}[Counterinstance]
Take $\dot I=1$, $P=0$, $D=0$, $k_B=T=1$. Every guard holds and the right-hand side is $0$, so
the conclusion would be $1\le0$.
\end{proof}
We gave a Lean proof of Theorem~\ref{thm:false} in session P5-T002, together with a lemma
\texttt{bound\_proves\_anything} showing that, taken as a hypothesis, the unrestricted bound
proves every proposition:
\begin{lstlisting}
theorem main : ViridisStatements.P5.UnrestrictedBoundFalse.Statement := by
  intro h
  have h1 := h 1 0 0 1 1 one_pos one_pos le_rfl le_rfl zero_le_one
  rw [bound_zero_power] at h1
  exact absurd h1 (by norm_num)
\end{lstlisting}
The witness is not degenerate: for any admissible $P,D,k_B,T$, the rate
$\mathrm{maxBitRate}(P,k_B,T)\cdot D+1$ violates the bound. According to the 2026-11-11 reading of
the corpus, no declaration used the axiom. That claim comes from reading the sources and is not
machine-checked. A well-posed successor would need a predicate that ties the rate to a modelled
physical system, together with separate empirical validation. Neither is attempted here.

\section{Finding 2: ``no finite compensation'' proves finite compensation}\label{sec:comp}

\subsection{The historical text}
The project README quotes \texttt{Non\_Substitutability\_Complete.lean} (SHA-256 prefix
\texttt{8b1194634\ldots bee273}), which introduces the theorem with the words ``there is no finite
power increase $\Delta P$ that can fully compensate''. Its formal text is:
\begin{lstlisting}
theorem no_finite_compensation (c : PhysicalConstants)
    {P D ΔD : ℝ} (hP : P > 0) (hD : D > 0) (hΔD : ΔD > 0) (hΔD' : ΔD < D) :
    let required_ΔP := P * ΔD / (D - ΔD)
    intelligence_bound c (P + required_ΔP) (D - ΔD) = intelligence_bound c P D
\end{lstlisting}
This is an \emph{existence} statement for a finite compensating increase. Since
\[
  P+\frac{P\,\Delta D}{D-\Delta D}=\frac{P(D-\Delta D)+P\,\Delta D}{D-\Delta D}
  =\frac{P\,D}{D-\Delta D},
\]
the theorem says that raising power to $PD/(D-\Delta D)$, which is finite whenever
$\Delta D<D$, restores the bound exactly after the loss. The name and the prose state the
opposite.

\subsection{The four facts}
Throughout, $\mathrm{cap}(P,D,C)=PD/C$ with $C>0$ and $P>0$.

\begin{theorem}[sealed statement \texttt{P5.CompensationUnderCeiling}]\label{thm:comp}
\leavevmode
\begin{enumerate}
\item \textbf{Interior.} If $0<\Delta D<D$, then
  $\mathrm{cap}\bigl(\tfrac{PD}{D-\Delta D},\,D-\Delta D,\,C\bigr)=\mathrm{cap}(P,D,C)$, and any
  $P'$ with $\mathrm{cap}(P',D-\Delta D,C)=\mathrm{cap}(P,D,C)$ equals $\tfrac{PD}{D-\Delta D}$.
\item \textbf{Ceiling.} If moreover $P\le P_{\max}$, then there is $P'\le P_{\max}$ with
  $\mathrm{cap}(P',D-\Delta D,C)=\mathrm{cap}(P,D,C)$ if and only if
  $\Delta D\le D\,(1-P/P_{\max})$.
\item \textbf{Divergence.} If $D>0$, then for every $M$ there is $\delta>0$ such that
  $\tfrac{PD}{D-\Delta D}>M$ whenever $D-\delta<\Delta D<D$.
\item \textbf{Complete loss.} If $D>0$, then for every $P'$,
  $\mathrm{cap}(P',0,C)=0<\mathrm{cap}(P,D,C)$.
\end{enumerate}
\end{theorem}

\begin{proof}[Proof sketch]
(1) $\frac{PD}{D-\Delta D}(D-\Delta D)=PD$ because $D-\Delta D>0$. Conversely, equal capacities
with the same $C>0$ give $P'(D-\Delta D)=PD$, so $P'=PD/(D-\Delta D)$.
(2) From $0<P\le P_{\max}$ we get $P_{\max}>0$. Multiplying by $P_{\max}$ shows
$\Delta D\le D(1-P/P_{\max})\iff PD\le P_{\max}(D-\Delta D)\iff
\tfrac{PD}{D-\Delta D}\le P_{\max}$. By the uniqueness in~(1), a compensating $P'\le P_{\max}$
exists exactly when the unique compensating power is at most $P_{\max}$.
(3) Take $\delta=PD/(|M|+1)$. If $0<D-\Delta D<\delta$, then
$\tfrac{PD}{D-\Delta D}>|M|+1>M$.
(4) $P'\cdot0/C=0$, and $PD/C>0$.
\end{proof}

We give a Lean proof of Theorem~\ref{thm:comp}. The automated checker run during the session
passed all five stages (denylist, compilation, statement match, axioms, kernel replay), and the
only axioms used are \texttt{propext}, \texttt{Classical.choice} and \texttt{Quot.sound}. The
key lemmas of the candidate file are:
\begin{lstlisting}
theorem interior (P D dD C : ℝ) (hC : 0 < C) (hP : 0 < P) (hdD : 0 < dD) (hlt : dD < D) :
    cap (P * D / (D - dD)) (D - dD) C = cap P D C ∧
    ∀ P' : ℝ, cap P' (D - dD) C = cap P D C → P' = P * D / (D - dD) := by
  have hpos : 0 < D - dD := by linarith
  refine ⟨?_, ?_⟩
  · unfold cap
    rw [div_mul_cancel₀ _ hpos.ne']
  · intro P' h
    unfold cap at h
    rw [div_left_inj' hC.ne'] at h
    rw [eq_div_iff hpos.ne', h]

theorem ceiling (P D dD C Pmax : ℝ) (hC : 0 < C) (hP : 0 < P) (hdD : 0 < dD) (hlt : dD < D)
    (hPmax : P ≤ Pmax) :
    (∃ P' : ℝ, P' ≤ Pmax ∧ cap P' (D - dD) C = cap P D C) ↔ dD ≤ D * (1 - P / Pmax) := by
  have hpos : 0 < D - dD := by linarith
  have hM : 0 < Pmax := lt_of_lt_of_le hP hPmax
  have key : dD ≤ D * (1 - P / Pmax) ↔ P * D ≤ Pmax * (D - dD) := by
    have : D * (1 - P / Pmax) = (Pmax * D - P * D) / Pmax := by
      field_simp
    rw [this, le_div_iff₀ hM]
    constructor <;> intro h <;> nlinarith
  rw [key]
  constructor
  · rintro ⟨P', hP', h⟩
    obtain ⟨-, huniq⟩ := interior P D dD C hC hP hdD hlt
    have he := huniq P' h
    rw [he, div_le_iff₀ hpos] at hP'
    exact hP'
  · intro h
    refine ⟨P * D / (D - dD), ?_, (interior P D dD C hC hP hdD hlt).1⟩
    rw [div_le_iff₀ hpos]
    exact h

theorem divergence (P D C : ℝ) (_hC : 0 < C) (hP : 0 < P) (hD : 0 < D) (M : ℝ) :
    ∃ δ : ℝ, 0 < δ ∧ ∀ dD : ℝ, D - δ < dD → dD < D → M < P * D / (D - dD) := by
  have hPD : 0 < P * D := mul_pos hP hD
  have hA : 0 < |M| + 1 := by positivity
  refine ⟨P * D / (|M| + 1), div_pos hPD hA, ?_⟩
  intro dD h1 h2
  have hpos : 0 < D - dD := by linarith
  have hsmall : D - dD < P * D / (|M| + 1) := by linarith
  rw [lt_div_iff₀ hA] at hsmall
  rw [lt_div_iff₀ hpos]
  have hM : M ≤ |M| := le_abs_self M
  nlinarith [abs_nonneg M]

theorem complete_loss (P D C : ℝ) (hC : 0 < C) (hP : 0 < P) (hD : 0 < D) (P' : ℝ) :
    cap P' 0 C = 0 ∧ cap P' 0 C < cap P D C := by
  have h0 : cap P' 0 C = 0 := by simp [cap]
  refine ⟨h0, ?_⟩
  rw [h0]
  unfold cap
  positivity
\end{lstlisting}
The theorem \texttt{main} is the conjunction of the four lemmas, each applied to the binders of
the sealed statement.

\subsection{Worked numerical examples}
Take $C=1$, $P=1$ and $D=1$ as the baseline, so that $\mathrm{cap}(P,D,C)=1$.

\begin{example}[Interior]
Lose half of $D$: $\Delta D=\tfrac12$. The compensating power is
$P'=\frac{1\cdot1}{1-1/2}=2$, and $\mathrm{cap}(2,\tfrac12,1)=1$. Power doubles to offset the
halving of $D$, and no other power restores the capacity. A smaller loss, $\Delta D=0.2$, needs
$P'=1/0.8=1.25$.
\end{example}

\begin{example}[Ceiling]
Keep $\Delta D=\tfrac12$. With $P_{\max}=2$ the condition reads
$\tfrac12\le1\cdot(1-\tfrac12)=\tfrac12$, which holds with equality: the required power $2$ is
exactly at the ceiling. With $P_{\max}=\tfrac32$ the threshold is $1\cdot(1-\tfrac23)=\tfrac13$
and $\tfrac12>\tfrac13$, so no admissible power compensates, matching $2>\tfrac32$. Under this
ceiling, every loss $\Delta D\le\tfrac13$ is compensable and every larger loss is not.
\end{example}

\begin{example}[Divergence]
The required power is $P'=1/(1-\Delta D)$: it is $10$ at $\Delta D=0.9$, $100$ at $0.99$ and
$1000$ at $0.999$. For the bound $M=1000$, the proof's choice $\delta=PD/(|M|+1)=1/1001$ gives
$P'>1001>1000$ whenever $1-\tfrac1{1001}<\Delta D<1$.
\end{example}

\begin{example}[Complete loss]
At $D=0$, $\mathrm{cap}(1000,0,1)=0<1=\mathrm{cap}(1,1,1)$, and the same holds with $1000$
replaced by any real power.
\end{example}

\begin{table}[h]
\centering
\begin{tabular}{@{}lccc@{}}
\toprule
Loss $\Delta D$ (with $P=D=C=1$) & Required $P'$ & Feasible if $P_{\max}=2$ & Feasible if $P_{\max}=\tfrac32$\\
\midrule
$0.2$ & $1.25$ & yes & yes\\
$1/3$ & $1.5$ & yes & yes (equality)\\
$0.5$ & $2$ & yes (equality) & no\\
$0.9$ & $10$ & no & no\\
$0.999$ & $1000$ & no & no\\
$1$ (complete) & none & no & no\\
\bottomrule
\end{tabular}
\caption{Compensating power for a range of losses, and whether it fits under two ceilings. Values
follow from $P'=PD/(D-\Delta D)$ and the ceiling condition of Theorem~\ref{thm:comp}(2).}
\end{table}

\subsection{Negative controls}
Both hypotheses of part~(2) besides $C>0$ and $P>0$ are needed. Without $P\le P_{\max}$: take
$P_{\max}=-1$, $P=D=C=1$, $\Delta D=\tfrac12$. The right side
$\tfrac12\le1\cdot(1-1/(-1))=2$ holds, but no $P'\le-1$ satisfies $P'\cdot\tfrac12=1$. Without
$\Delta D<D$: take $P=D=C=P_{\max}=1$ and $\Delta D=2$. Then $P'=-1\le P_{\max}$ restores the
capacity (both are $1$), but $2\le1\cdot(1-1)=0$ fails. Both controls are checked in the sealed
module.

\subsection{The economic reading}
Along a level set $PD=\text{const}$ the two inputs trade off one for one in proportional terms:
losing a fraction $\Delta D/D$ of diversity is offset by multiplying power by $D/(D-\Delta D)$.
Three things should be kept apart.
\begin{itemize}
\item \emph{Essential.} $D$ is essential: at $D=0$ capacity is $0$ for every power
  (part~4). Nothing replaces $D$ at zero.
\item \emph{Not non-substitutable in the interior.} For any partial loss, a finite power restores
  capacity exactly (part~1). In this algebra, $D$ and $P$ are substitutes away from the boundary.
  The historical theorem itself proves this.
\item \emph{Ceiling.} What makes a partial loss uncompensable is a finite power ceiling
  (part~2). The tolerable loss $D(1-P/P_{\max})$ shrinks to zero as current power approaches the
  ceiling. Divergence (part~3) explains why every finite ceiling eventually binds: losses close
  enough to total need more power than any fixed bound.
\end{itemize}
Any argument that losses of $D$ cannot be compensated must therefore rest on a ceiling on power
(or on some other constraint outside this algebra), not on the capacity formula alone.

\section{What is not shown}
These results concern the formal text and the algebra $\mathrm{cap}=PD/C$. They do \emph{not}
show:
\begin{itemize}
\item any behaviour of any agent, for example that an agent would or should preserve
  biodiversity. A capacity bound does not say it is attained, that more capacity serves an
  agent's goals, or that preservation is the best action once costs and alternatives are counted;
\item any measured value of $D$, or that $D$ as used here corresponds to any measurable ecological
  quantity;
\item any ecological claim, including any claim about how much diversity loss is tolerable in
  practice, or about the size of any real power ceiling;
\item anything about the Intelligence Bound as a physical conjecture.
\end{itemize}

\section{Limitations}
Both results are elementary, and they are corrections, not new theorems. The sealed statements
restate the historical definitions instead of importing the historical files, so their fidelity
to the originals rests on the quoted text and the file hashes. The statement that no corpus
declaration depended on the false axiom comes from reading the sources and is not machine-checked.
The capacity model multiplies $D$ in by fiat. A different functional form, for instance one with
diminishing returns to power, would change the compensating power and the ceiling condition.
These results do not cover such forms.

\section{Reproducibility}
The proofs are \texttt{proofs/P5.UnrestrictedBoundFalse.lean} (session P5-T002) and
\texttt{proofs/P5.CompensationUnderCeiling.lean} (this session). Each imports only Mathlib and
its sealed module: \texttt{ViridisStatements.P5.UnrestrictedBoundFalse} (SHA-256
\path{3b83eb37905660e3db1cc68c54e476f4ac5c450b245cec4675195f9cc2e0a95d}) and
\texttt{ViridisStatements.P5.CompensationUnderCeiling} (SHA-256
\path{ff09ca9eb882869dcbff84346866ae7b501114797ce2a5f97865260552ab763b}). They compile with
toolchain \texttt{leanprover/lean4:v4.28.0} and Mathlib commit
\path{8f9d9cff6bd728b17a24e163c9402775d9e6a365}, and can be checked with
\texttt{viridis-lean-check proofs/<ID>.lean <ID>}. The worked examples are hand arithmetic and
are also checked as sanity examples in the sealed module. No data or external code is used.

\section{Changelog}
\begin{itemize}
\item 2026-11-13 (P5-T001): first version. Proposed both statements with sanity checks and
negative controls.
\item 2026-11-14 (P5-T002): added a Lean proof of the sealed \texttt{P5.UnrestrictedBoundFalse}
and an explosion lemma. The second correction remained a conjecture.
\item 2026-11-19 (P5-T003): added a Lean proof of the sealed \texttt{P5.CompensationUnderCeiling},
replacing the conjecture. Quoted \texttt{no\_finite\_compensation} and showed that it proves
finite compensation. Set out the four facts with worked examples, a table and negative controls,
and added the economic reading (essential but substitutable in the interior; a ceiling makes
losses uncompensable) and a section on what is not shown. Condensed Finding~1 and retitled the
note to cover both corrections.
\end{itemize}

\bibliographystyle{plainnat}
\bibliography{refs}
\end{document}
